\documentclass[a4paper,12pt,fleqn]{article}%fleqn pour aligner les equations à gauche
\usepackage{Eval}
\usepackage[toc,page]{appendix}


\newcommand{\chnum}{Chapitres 3 et 12}
\newcommand{\chtitle}{Cinétique chimique et mouvement dans un champ uniforme}
\newcommand{\dtype}{DS 4}

\begin{document}
\maketitle
\section{Détecteur ionique de fumées}
\begin{figure}[h!]
    \begin{minipage}{.48\linewidth}
    Le principe de ce détecteur de fumées repose sur l’ionisation de l’air par des particules $\alpha$. En l’absence de fumées, ces particules arrachent des électrons aux molécules de dioxygène  et de diazote présentes dans la chambre à ionisation. Pour le dioxygène, l’ionisation nécessite un apport d’énergie de 12 eV par molécule.\\
    Les ions et les électrons formés par l’ionisation de l’air sont soumis à un champ électrique uniforme entre deux plaques. Un courant électrique de faible intensité apparaît alors dans le circuit électrique (Figure~\ref{Fig1}).\\
    Lorsque la fumée pénètre dans la chambre à ionisation, une partie des électrons et des ions issus de l’ionisation se fixe aux poussières de fumées. La baisse de l’intensité électrique qui en résulte déclenche un avertisseur sonore.
\end{minipage}
\begin{minipage}{.5\linewidth}
    \includegraphics[width=\linewidth]{Ressources/TSPE_DS4_24b_Fig1.png}
    \caption{Illustration du principe d'une chambre à ionidation}
    \label{Fig1}
\end{minipage}
\end{figure}
\begin{figure}[h!]
    \begin{minipage}{.48\linewidth}
    On s’intéresse au mouvement d’une particule $\alpha$ arrivant dans la chambre à ionisation en l’absence de fumée. Cette particule arrive en un point $O$ avec un vecteur vitesse initiale parallèle aux plaques $C$ et $D$ du condensateur plan (voir Figure~\ref{Fig2}).\\
    Une tension constante $U = \SI{9,0}{V}$ est appliquée entre les deux plaques $C$ et $D$. La valeur de la vitesse initiale $v_O$ est égale à \SI{1,6E7}{\meter\per\second}.\\
    On étudie le mouvement de la particule $\alpha$ dans le référentiel terrestre supposé galiléen. À l’instant $t = 0$, la particule $\alpha$ est au point O.\\
    Lors de cette étude, on négligera les éventuelles collisions avec les molécules de l’air ainsi que la valeur du poids de la particule $\alpha$ devant la valeur de la force électrique $\vv{F_e}$ subie par cette particule.
\end{minipage}
\begin{minipage}{.5\linewidth}
    \includegraphics[width=\linewidth]{Ressources/TSPE_DS4_24b_Fig2.png}
    \caption{Condensateur plan utilisé dans le détecteur de fumée}
    \label{Fig2}
\end{minipage}
\end{figure}

\textbf{Données :}
\begin{itemize}
    \item 1 électronvolt (eV) = \SI{1,5E-19}{J}
    \item charge élémentaire : $e=\SI{1,6E-19}{C}$
    \item Pour un condensateur plan, le champ électronique $E$ est relié à la tension $u$ à à la distance $d$ qui sépare les deux plaques par la relation $E=\dfrac{U}{d}$
    \item Charge de la particule $\alpha$ : $q_{\alpha} = + 2 e$
    \item Masse d'une particule $\alpha$ : $m_{\alpha} = \SI{6,64E-27}{kg}$
    \item Intensité du champ de pesanteur terrestre : $g = \SI{9,81}{\meter\per\second}$
\end{itemize}


\begin{enumerate}
    \item Vérifier par le calcul que l’hypothèse concernant le poids de la particule $\alpha$ est justifiée.
    \item Reproduire sur la copie le schéma de la figure~\ref{Fig1} puis y représenter le champ électrostatique $\vv{E}$ et la force électrique que subit la particule $\alpha$ au point $O$. Justifier.
    \item Établir que les équations horaires du mouvement de la particule $\alpha$ sont : $$x(t)=v_0 \cdot t$$
    $$y(t)= \dfrac{e\cdot U}{m_{\alpha}\cdot d}\cdot t^2$$
    \item Déterminer la valeur de la coordonnée $y_L$ de la particule lorsqu’elle a parcouru une distance suivant l’axe $Ox$ égale à $L = \SI{4,0}{cm}$. Expliquer pourquoi le mouvement de cette particule peut être considérée comme rectiligne dans la chambre d’ionisation.
    \item Montrer que l’énergie cinétique initiale des particules $\alpha$ est suffisante pour ioniser des molécules de dioxygène.
\end{enumerate}


\section{Étude cinétique d’une réaction d’oxydoréduction}
On se propose d’étudier la cinétique de la transformation lente de décomposition de l’eau oxygénée par les ions iodure en présence d’acide sulfurique, transformation considérée comme totale.\\
L’équation de la réaction qui modélise la transformation d’oxydoréduction s’écrit :
$$\ce{H2O2(aq) + 2 I-(aq) + 2 H+(aq) <=> I2(aq) + 2H2O(l)}$$

La solution de diiode formée étant colorée, la transformation est suivie par spectrophotométrie, méthode qui consiste  à mesurer l'absorbance A de la solution, grandeur proportionnelle à la concentration en diiode.

	\subsection{Étude théorique de la réaction}
	\begin{enumerate}
		\item Identifier,  dans l'équation de la réaction étudiée, les deux couples d'oxydoréduction mis en jeu en écrivant leurs demi-équations correspondantes.
		Justifier l’écriture de l’équation de la réaction.
	\end{enumerate}
	\subsection{Suivi de la réaction}
	À la date $t$= \SI{0}{s}, on mélange un volume $V_1$ = \SI{20,0}{mL} d'une solution d'iodure de potassium (\ce{K+} ; \ce{I-}) de concentration $C_1$ = \SI{2,0}{mL} d'eau oxygénée \ce{H2O2} de concentration $C_2$ = \SI{0,10}{\mole\per\liter}.\\
	On remplit une cuve de spectrophotométrie et on relève les valeurs de l'absorbance au cours du temps. On détermine alors, grâce à la loi de Beer-Lambert, la concentration en quantité de matière  [\ce{I2}] du diiode formé :
    \begin{center}
        \begin{tabular}{|c|c|c|c|c|c|c|c|}
		\hline t (s)& 0 & 126 & 434 & 682 & 930 & 1178 & 1420\\
		\hline [\ce{I2}] (\unit{\milli\mole\per\liter}) & 0,00 & 1,74 & 4,06 & 5,16 & 5,84 & 6,26 & 6,53\\
		\hline
	\end{tabular}
    \end{center}
	
	
	\begin{enumerate}
		\item Justifier l'utilisation du spectrophotomètre pour le suivi cinétique de la réaction.
		\item Rappeler la relation entre l'absorbance $A$ et la concentration en quantité de matière [\ce{I2}] de l'espèce colorée.
		\item Le mélange initial est-il dans les proportions sot\oe chiométriques ? Justifier la réponse par des calculs.
		\item Compléter le tableau d'avancement de la transformation et justifier les calculs de $n_1$ et $n_2$.
        \end{enumerate}
		\noindent \begin{tabular}{|>{\centering}m{3cm}|>{\centering}m{2.5cm}|>{\centering}m{2cm}|>{\centering}m{2cm}|>{\centering}m{2cm}||>{\centering}m{2cm}|>{\centering\arraybackslash}m{2cm}|}
			\hline
			\multicolumn{2}{|c|}{Équation chimique} & \multicolumn{5}{|c|}{\ce{H2O2(aq) + 2 I-(aq) + 2 H+(aq) <=> I2(aq) + 2H2O(l)}}\\
			\hline État du système & Avancement (\unit{\mole}) & \multicolumn{5}{|c|}{Quantité de matière (\unit{\mole})}\\
			\hline État initial & 0 & $n_2$ = \num{2,0e-4} & $n_1$ = \num{20e-4} & Excès & Beaucoup & \\
			\hline État en cours de transformation & $x$ & & & Excès & Beaucoup & \\
			\hline État final & $x_{max}$ & & & Excès & Beaucoup & \\
			\hline
		\end{tabular}
        \begin{enumerate}[resume]
		\item Déterminer l'avancement maximal. En déduire la valeur théorique de la concentration en diiode formé lorsque la transformation est terminée. Cette concentration est-elle en accord avec le tableau de mesures ?
	\end{enumerate}
	
	\subsection{Exploitation des résultats}
	À l'aide d'un tableur, on calcule la concentration et la vitesse de disparition de l'eau oxygénée au cours du temps. Une capture d'écran est donnée ci-dessous.\\
	\begin{center}\includegraphics[width = .5\linewidth]{Ressources/TSPE_C10_Eval_Fig8.png}\end{center}
	\begin{enumerate}
		\item Donner la définition du temps de demi-réaction dans le cas étudié, puis déterminer graphiquement $t_{1/2}$ à l'aide de la courbe de l'évolution de la concentration en eau oxygénée au cours du temps ci-après.\\
		\begin{center}\includegraphics[width = .5\linewidth]{Ressources/TSPE_C10_Eval_Fig6.png}\end{center}
		\item  Retrouver par un calcul, la vitesse de disparition de \ce{H2O2} à $t$ = \SI{434}{s}.\\
		On trace ensuite la courbe donnant la vitesse de disparition $v_D$ de l'eau oxygénée en fonction de sa concentration.\\
		\begin{center}\includegraphics[width = .5\linewidth]{Ressources/TSPE_C10_Eval_Fig7.png}\end{center}
		\item Expliquer pourquoi cette transformation chimique suit une loi de vitesse d'ordre 1.
		\item La modélisation de la concentration en eau oxygénée au cours du temps est donnée par la relation
        $$[\ce{H2O2}](t) = [\ce{H2O2}]_0 \cdot  e^{-kt}$$
		À l'aide de cette relation et du tableau de mesures, déterminer [\ce{H2O2}]$_0$, concentration initiale en eau oxygénée.
		\item Retrouver cette valeur par le calcul.
		\item Déterminer $k$ la constante de vitesse en utilisant la droite donnant la vitesse de disparition $v_D$ de l'eau oxygénée en fonction de sa concentration.
	\end{enumerate}
\end{document}
